\section{Theoretischer Hintergrund}

\subsection{Industrieroboter}
Ein Industrieroboter, ist gemäß DIN EN ISO 10 218-1 ein frei Programmierbarer Mehrzweck-Manipulator, der in 3 oder mehr Achsen Programmierbar ist.
Er ist universell einsetzbar, und wird häufig für Schweiß-, -Lackierarbeiten, Palletieraufgaben und viele weitere Fertigungs- und Handhabeaufgaben verwendet.
Obwohl ein Industrieroboter, im Gegensatz zu mobilen Robotern sich nicht frei durch den Raum bewegen kann, kann er statt an festen Orten auch beweglich angeordnet werden \doublecite[17-18, 1-2]{Industrieroboter, weber-Industrieroboter}.


Im Gegensatz zu mobilen Robotern, die sich im Raum fortbewegen können,
bleibt ein Manipulator an einem festen Ort und führt Bewegungen über seine Gelenke und Glieder aus.
Typischerweise besteht ein solcher Roboter aus mehreren starren Segmenten (sogenannten Gliedern),
die durch Gelenke miteinander verbunden sind.
Diese Struktur erlaubt es dem Roboter, präzise Bewegungen in mehreren Freiheitsgraden durchzuführen.



Ihre Bewegungen werden entweder durch direkte Programmierung, durch sogenannte Teach-In-Verfahren oder durch komplexe Steuerungsalgorithmen bestimmt, die z.B.
Pfadplanung und inverse Kinematik umfassen.

Ein klassisches Beispiel für einen Manipulator ist der Industrieroboterarm.
Er imitiert in gewissem Maße den menschlichen Arm: Er hat typischerweise eine Basis (entsprechend der Schulter),
einen Arm (Oberarm und Unterarm), ein Handgelenk sowie gegebenenfalls ein Greifwerkzeug (Endeffektor) an der Spitze.
Die Steuerung solcher Systeme erfordert ein Verständnis für Robotikgrundlagen wie Kinematik, Dynamik, Sensorik und Steuerungstechnik.

\subsection{Vorwärtskinematik}
Die Vorwärtskinematik (auch direkte Kinematik genannt) beschreibt in der Robotik den Zusammenhang zwischen den Gelenkwinkeln eines Roboters und der daraus resultierenden Position und Orientierung des Endeffektors im Raum.
Sie ist ein zentrales Konzept für die Steuerung und Simulation von Manipulatoren, da sie es ermöglicht, den Zustand des Roboterarms basierend auf den aktuellen Gelenkstellungen zu berechnen.

Mathematisch wird die Vorwärtskinematik durch Transformationen beschrieben, die die Position und Ausrichtung eines Glieds relativ zu einem vorhergehenden Glied bestimmen.
Dazu werden in der Regel Homogene Transformationsmatrizen verwendet, die sowohl Rotation als auch Translation in einer einzigen Matrix kombinieren.
Durch die sukzessive Multiplikation dieser Matrizen entlang der kinematischen Kette – von der Basis bis zum Endeffektor – ergibt sich die absolute Pose des Endeffektors im kartesischen Raum.

Ein verbreitetes Modell zur systematischen Beschreibung der Vorwärtskinematik ist die Denavit-Hartenberg-Notation.
Dieses Verfahren erlaubt es, die Transformationen zwischen benachbarten Gliedern mit nur vier Parametern pro Gelenk kompakt darzustellen.

Die Berechnung der Vorwärtskinematik ist deterministisch und liefert für jeden Satz an Gelenkwerten genau eine Lösung.
Dies unterscheidet sie von der inversen Kinematik, bei der für eine bestimmte Zielposition mehrere, eine oder gar keine Lösung existieren kann.
In vielen Robotiksystemen ist die Vorwärtskinematik Bestandteil der grundlegenden Steuerlogik und wird auch für Visualisierung und Kollisionserkennung eingesetzt.




\subsection{Denavit-Hartenberg-Notation}
Die Denavit-Hartenberg-Notation (kurz: DH-Notation) ist ein standardisiertes Verfahren zur Modellierung der Kinematik serieller Roboterarme.
Sie wurde 1955 von Jacques Denavit und Richard Hartenberg eingeführt und ist bis heute eines der am weitesten verbreiteten Werkzeuge zur Beschreibung von Roboterkinematiken.
Ziel der Notation ist es, die Position und Orientierung jedes Gliedes relativ zum vorhergehenden Glied mit einem minimalen Satz von Parametern darzustellen.

Bei der DH-Notation wird jedem Gelenk ein eigenes Koordinatensystem zugewiesen.
Die Transformation von einem Glied zum nächsten wird durch eine sogenannte homogene Transformationsmatrix beschrieben, die vier aufeinanderfolgende Basisoperationen zusammenfasst:
zwei Rotationen und zwei Translationen. Diese vier DH-Parameter sind:

  $\theta_i$: der Gelenkwinkel, Rotation um die z-Achse

  $d_i$: die Verschiebung entland der z-Achse

  $a_i$: die Länge des Arms, also die verschiebung entlang der x-Achse

  $\alpha_i$: der Verwindungswinkel, Rotation um die x-Achse

  Mit diesen Parametern lässt sich für jedes Glied eine 4×4-Transformationsmatrix aufstellen, die die Position und Orientierung relativ zum vorherigen Koordinatensystem beschreibt.
  Die Gesamttransformation vom Basis-Koordinatensystem zum Endeffektor ergibt sich durch sukzessive Multiplikation aller einzelnen Transformationsmatrizen entlang der kinematischen Kette.

Ein Vorteil der DH-Notation liegt in ihrer systematischen und einheitlichen Struktur,
was insbesondere bei der Programmierung von Robotern und bei der Analyse komplexer Kinematiken hilfreich ist.
Allerdings kann die Anwendung bei bestimmten Robotertypen – z.B. bei parallelen Strukturen oder redundanten Systemen – unübersichtlich oder eingeschränkt sein.
In solchen Fällen kommen häufig alternative Modellierungsverfahren wie die modifizierte DH-Notation oder Produkt-von-Exponentialen-Darstellungen zum Einsatz.



\subsection{Rückwärtstransformation}
Die Rückwärtstransformation – häufig auch als inverse Transformation bezeichnet – beschreibt den rechnerischen Vorgang,
bei dem aus einer bekannten Pose des Endeffektors (Position und Orientierung im Raum) auf die zugehörigen Gelenkparameter eines Roboters geschlossen wird.
Sie ist das Gegenstück zur Vorwärtskinematik und stellt eine zentrale Herausforderung in der Robotik dar.

Im Gegensatz zur Vorwärtskinematik, bei der es für eine bestimmte Konfiguration der Gelenke genau eine Lösung für die Endeffektorposition gibt,
ist die Rückwärtstransformation in der Regel nicht eindeutig lösbar. Für viele Aufgabenstellungen existieren mehrere mögliche Gelenkstellungen,
die dieselbe Endeffektorposition erreichen können – oder in manchen Fällen auch keine.
Diese Mehrdeutigkeit entsteht insbesondere bei Robotern mit sechs oder mehr Freiheitsgraden.

Die mathematische Lösung der Rückwärtstransformation kann analytisch oder numerisch erfolgen.

Analytische Methoden liefern geschlossene Formeln für die Gelenkwinkel, sind aber nur für bestimmte Roboterstrukturen möglich (z.B. bei standardisierten 6-DOF-Industrierobotern mit spezifischer Gelenkanordnung).

Numerische Methoden wie das Newton-Raphson-Verfahren oder Gradientenverfahren arbeiten iterativ und können auch für komplexere oder redundante Roboter eingesetzt werden.
Allerdings erfordern sie eine geeignete Startlösung und können je nach Problemstellung langsam oder instabil sein.

Ein häufig genutzter Ansatz zur Implementierung der Rückwärtstransformation ist die Kombination aus Vorwärtskinematik und Optimierung:
Ein Ziel im kartesischen Raum wird durch sukzessive Anpassung der Gelenkwinkel angesteuert, bis die berechnete Pose mit der Zielpose innerhalb einer akzeptablen Toleranz übereinstimmt.

Die Rückwärtstransformation spielt eine wesentliche Rolle bei der Trajektorienplanung, da in vielen Anwendungen die gewünschte Bewegung im kartesischen Raum beschrieben wird (z.B. entlang einer Linie oder Kurve),
während die Steuerung des Roboters auf Gelenkebene erfolgen muss.